\section{Introduction}\label{sec:introduction}

This paper studies a signed analogue of the type-A one-incidence layer.  In type A, the layer is defined by one hit on the identity diagonal and one hit on the reflected diagonal determined by the reversal permutation.  Here the ambient group is the signed permutation group $B_n$, and the signed layer records positive hits on the identity diagonal and on the reversal diagonal.  The reversal diagonal is essential: it is the geometric second diagonal that makes the object noncentral but still rigid enough to be controlled by the matching scheme of the reversal pairs.

The main theorem is deliberately scoped.  We do not attempt a full native $B_n$ Fourier fingerprint or a general Weyl-group theorem.  We prove the exact ranks of two ambient standard-channel aggregate operators on every nonempty diagonal row $\BtildeI_n(k,k)$.  The numbers that appear, $\lceil n/2\rceil$ and $\lceil n/2\rceil-1$, are classical multiplicities from the matching scheme attached to the reversal involution.  The object-specific contribution is that this noncentral signed-reversal layer collapses onto the corresponding matching-scheme components, and that the remaining scalars are strictly positive.  We then extend the two rank equalities to the type-D parity sublayer $\DtildeI_n(k,k)=\BtildeI_n(k,k)\cap D_n$, where $D_n=\ker\chi$ is the even-sign subgroup: a sign-character split writes the type-D aggregate as the average of the sign-blind and the sign-weighted aggregates, and the sign-weighted correction is controlled by a factorized generating function and, for the odd pair-standard channel, by a direct type-D trace identity.

There are two independent additions around the main theorem.  First, the positive-only row $X_n=\BtildeI_n(1,1)$ is separated from two natural chamber families: for $n\ge4$ it is neither a Reiner-style signed-poset linear-extension set nor an exact left-translate tiler of $B_n$.  Second, the appendix records finite $n=8$ native $B_n$ fingerprint evidence.  That evidence explains the known zero rows in one finite atlas slice, but it is not promoted to a classification theorem.

The proof of the rank theorem has three parts.  The even rows collapse directly by centralizer symmetry and a permanent-domination positivity lemma.  The odd rows reduce to two residual scalar sequences, called $H_m$ and $K_m$.  These residuals are identified with signed-rook moment formulae and then proved positive by exact order-three recurrences and ratio induction.  The computational material in the repository verifies the polynomial identities and finite edge rows using exact integer or rational arithmetic; the paper states the mathematical identities that those certificates check.

\paragraph{Structure.}
Section~\ref{sec:defs} fixes the signed-reversal layer and states the main theorem.  Section~\ref{sec:separator} proves the positive-only separator theorem.  Sections~\ref{sec:background}--\ref{sec:assembly} prove the rank theorem.  Section~\ref{sec:typed} extends it to the type-D parity sublayer.  Section~\ref{sec:boundaries} records the claim boundary.  The appendices give the exact replay routing and the finite native $B_n$ fingerprint boundary.
